\section{sPAR}
\begin{itemize}
    \item If the parameters support even the base case of $\eta = 2$ users, then the scheme can support a nearly arbitrary amount of users with the same parameters
    \item Noise is a big problem due to the internal product
    \item Getting enough noise headroom at $\lambda = 128$ requires very large parameters, making the protocol much slower than anticipated
\end{itemize}

\subsubsection{Noise}
% \definecolor{gaussblue}{RGB}{31,119,180}
% \definecolor{hpsgreen}{RGB}{44,160,44}
% \definecolor{canonorange}{RGB}{255,165,0}
% \definecolor{maxred}{RGB}{255,36,0}
% \definecolor{hpspurple}{RGB}{148,103,189}

% ---------- Compute fits in linear space ----------
% \pgfplotstableread[col sep=comma]{Data/Noise/write-N2.csv}\bvsmall
% \pgfplotstablecreatecol[linear regression={x=n, y=noise}]{reg}{\bvsmall}
% \edef\smallSlope{\pgfplotstableregressiona}
% \edef\smallIntercept{\pgfplotstableregressionb}

\pgfplotstableread[col sep=comma]{Data/Noise/write-N32.csv}\bvlarge
\pgfplotstablecreatecol[linear regression={x=n, y=noise}]{reg}{\bvlarge}
\edef\largeSlope{\pgfplotstableregressiona}
\edef\largeIntercept{\pgfplotstableregressionb}

\begin{figure}[H]
    \centering
    \caption{Noise in $L$ after full round.}
    \label{fig:spar:plot-1}
    \begin{tikzpicture}
    \begin{semilogyaxis}[
        noiseaxis,
        log basis y=2,
        xmin=1.1, xmax=33,
        xlabel={$\eta$},
        ylabel={Error $\|\epsilon\|_\infty$},
        legend pos=south east,
        set layers,
        width=\textwidth,
        height=0.55\textwidth,
    ]
        % ---------- Scatter plots ----------
        % \addplot[only marks, mark=*, mark size=4pt, color=gaussblue]
        %     table[col sep=comma, x expr=\thisrow{n}+1, y=noise] {Data/Noise/write-N2.csv};
        \addplot[only marks, mark=*, mark size=1pt, color=hpsgreen]
            table[col sep=comma, x expr=\thisrow{n}+1, y=noise] {Data/Noise/write-N32.csv};
            
        % ---------- Regression lines ----------
        % \addplot[thick, solid, dotted, on layer=axis foreground, color=canonorange, domain=1:16, samples=16, forget plot]
        %     {\smallSlope*x + \smallIntercept};
        % \addplot[thick, solid, dotted, on layer=axis foreground, color=hpspurple, domain=1:16, samples=16, forget plot]
        %     {\largeSlope*x + \largeIntercept};
            
        \legend{$\ell = 2$} % \legend{$\eta = 2$, $\eta = 16$}
    \end{semilogyaxis}
    \end{tikzpicture}
\end{figure}


% \newpage

% \begin{tabular}{lc} 
%     \toprule
%     \textbf{n} & \textbf{noise} \\ % & \textbf{Role} \\ % Manual Table Header
%     \midrule
    
%     % Reads CSV, maps column headers to macros, and loops through rows
%     \csvreader[head to column names]{Data/Noise/ExternalProd/bv.csv}{}
%     {\n & \noise \\} 
    
%     % \bottomrule
% \end{tabular}

\begin{itemize}
    \item Linear regression: $\largeSlope x + \largeIntercept$
    \item Noise ceiling: $180$ bits
    \item Assuming same growth: $\largeSlope \eta + \largeIntercept \leq 2^{180} \implies \eta \leq 1.03875\cdot10^{12}$ 
    \item Noise growth is probably larger as $\eta$ increases; have to verify
\end{itemize}


%%
%% Client::Encrypt, Server::Write, Client::Decrypt, Server::Decrypt
\subsection{Performance}
% TODO: Load from JSON again
\begin{table}[H]
    \caption{One round of the full sPAR protocol, not adjusted to per user, 6 threads}
    \label{tab:bench:full}
    \centering
    \begin{tabular}{cccccc}
        \toprule
        $\mathbf{\eta}$ & \textbf{Total time} & \textbf{Encrypt} & \textbf{Write} & \textbf{Partial Dec.} & \textbf{Server Dec.} \\
        \midrule
        2 & $78.586$ s  & $826.96$ ms   & $76.61$ s   & $201.83$ ms  & $5.12$ ms \\
        4 & $315.636$ s & $3049.33$ ms  & $308.01$ s  & $546.20$ ms  & $10.55$ ms \\
        8 & $1258.86$ s & $13562.60$ ms & $1223.85$ s & $1631.47$ ms & $26.61$ ms \\
        % \JSONParseArrayMapInline{\fullData}{benchmarks}{
        %   $\nvalue{#1}$ &
        %   \num{\jsonms{\fullData}{benchmarks[#1].cpu_time}} & 
        %   \num{\jsonmspern{\fullData}{benchmarks[#1].encrypt}{\nvalue{#1}}} &
        %   \num{\jsonmspern{\fullData}{benchmarks[#1].write}{\nvalue{#1}}} &
        %   \num{\jsonmspern{\fullData}{benchmarks[#1].partial}{\nvalue{#1}}} &
        %   \num{\jsonms{\fullData}{benchmarks[#1].final}} \\
        % }
        \bottomrule
    \end{tabular}
\end{table}

\begin{itemize}
    \item If write can use SIMD batching the write loop can be divided by up to $N \implies 4.675$ ms, $18.798$ ms, $74.69$ ms per round, respectively, but the other functions must also be scaled up.. % works in some scenarios
    \item \textcolor{red}{TODO: run single-threaded, adjust to $\eta$ i.e. divide by $\eta$ for encrypt, write, partial decrypt}
\end{itemize}

% Check gain vs DC-net -> throughput